\section{Certified traces, cyclic demands, and relational paths}
\label{sec:traces}

\subsection{Deterministic trace reflection}

Let a fixed state machine have primitive-recursive initial state
\(s_0(\vec q)\), step
\(s_{t+1}=F(t,s_t,\vec q)\), and a quantifier-free invariant
\(I(t,s,\vec q)\).  The state sequence is defined by \code{recDef}.

\begin{theorem}[Trace invariant]
\label{thm:trace}
Assume object derivations
\begin{align*}
 &\Gamma\turnstilep I(0,s_0(\vec q),\vec q),\\
 &I(t,s_t,\vec q),\Gamma^{\shift}
     \turnstilep I(t+1,F(t,s_t,\vec q),\vec q).
\end{align*}
Then iPRA proves \(I(T,s_T,\vec q)\) for every term \(T\).
\end{theorem}

\begin{proof}
Use quantifier-free induction on the formula
\(I(t,s_t,\vec q)\).  The zero and successor equations of \code{recDef}
rewrite the two induction cases to the supplied derivations.  Instantiate the
open result at \(T\).
\end{proof}

A human invariant containing bounded quantifiers is first compiled to a bit
\(i(t,s,\vec q)\); induction is applied to the equation \(i=1\).  Open
reflection recovers the human statement at the terminal state.

\subsection{Query-first precision scheduling}

An analytic or geometric computation should not carry an entire infinite name
through one loop.  Fix the terminal query \(q\) first, propagate the finite
input precision and resource requirements backward, and specialize every node
to those demands.  The forward trace then manipulates only finite state.  The
final theorem may quantify over \(q\) outside the trace.

This order is logically important.  It turns ``for every precision there is a
finite computation'' into one open theorem parameterized by precision, without
placing a universal quantifier inside the induction invariant.

\subsection{Finite monotone cyclic demand closure}

A reverse demand graph may have cycles.  Let
\[
  D=\prod_{j<m}\{0,\ldots,C_j\}
\]
with componentwise order, and let \(F:D\to D\) be monotone.  Define the
inflationary step
\[
  T(d)=d\vee F(d)
\]
with componentwise maximum and clipping at capacities.  Starting from a base
demand \(d_0\), put \(d_{r+1}=T(d_r)\).

\begin{theorem}[Capped demand stabilization]
\label{thm:demand}
Let \(H=\sum_{j<m}C_j\).  Then \(d_H\) is a prefixed point of \(F\).  It is the
least prefixed point above \(d_0\).
\end{theorem}

\begin{proof}
The sequence is componentwise nondecreasing.  Define the rank
\(\rho(d)=\sum_j d_j\).  If \(d_{r+1}\neq d_r\), at least one coordinate
strictly increases, so \(\rho(d_{r+1})\ge\rho(d_r)+1\).  The rank is at most
\(H\).  Therefore there can be at most \(H-\rho(d_0)\) strict steps; by stage
\(H\), \(T(d_H)=d_H\), hence \(F(d_H)\le d_H\).

For leastness, let \(p\ge d_0\) satisfy \(F(p)\le p\).  Induct on \(r\).  If
\(d_r\le p\), monotonicity gives \(F(d_r)\le F(p)\le p\), so
\(d_{r+1}=d_r\vee F(d_r)\le p\).  Thus \(d_H\le p\).
\end{proof}

For fixed dimension, all coordinates, comparisons, maxima, and the rank are PR
terms.  The two inductions use quantifier-free vector equations and order
formulas.  An overflow bit records whether the chosen caps were adequate; its
vanishing is either proved or retained as a residual.

\subsection{Relational paths and backward viability}

Pointwise existence at successive stages does not imply that independently
chosen witnesses form a compatible path.  Let finite layers
\(X_0,\ldots,X_T\) have reflected validity predicates and reflected edge
relations \(E_t(x,y)\).  Define viability backward:
\[
\begin{aligned}
  \mathsf{viable}_T(x)&=\mathsf{accept}(x),\\
  \mathsf{viable}_t(x)&=1
   \quad\Longleftrightarrow\quad
   \exists y<B_{t+1}\,[E_t(x,y)\wedge\mathsf{viable}_{t+1}(y)=1].
\end{aligned}
\]
Each existential is a bounded reflected search.

\begin{theorem}[Least viable path]
If an initial viable node exists, bounded least search at each forward stage
constructs a whole path whose edges and terminal acceptance are verified.
Conversely, viability zero implies that no accepted path begins there.
\end{theorem}

\begin{proof}
Backward induction establishes the equivalence between the viability bit and
existence of an accepted suffix.  In the forward pass, the current node's
viability supplies an accepted successor below the finite bound; choose the
least one by \cref{thm:least-code}.  Continue for the fixed horizon.  The
terminal viability equation is exactly acceptance.  The no-path direction
opens any purported path and propagates its suffix backward, contradicting
viability zero.
\end{proof}

This theorem resolves hidden finite choice.  A greedy locally valid choice is
never assumed compatible; the whole finite interaction is verified before any
representative path is named.
